\begin{recipe}{Eiwitpannenkoekjes met banaan}
\label{recipe:eiwitpannenkoekjes-banaan}
\kicker[Dessert,Zoet,Kinderen]{Dessert · zoet · kinderen}
\index[register]{Eiwitpannenkoekjes met banaan}
\index[register]{Pannenkoeken}
\index[register]{Eiwit}
\index[register]{Banaan}

\meta{10 minuten \dvd Voor 1 portie \dvd Snel}

\marginimage{images/eiwitpannenkoekjes-banaan}
\lettrine{D}{eze} pannenkoekjes zijn ontstaan omdat we nog drie eiwitten over
hadden van de pasta carbonara (p.~\pageref{recipe:pasta-carbonara}) en er
toevallig een overrijpe banaan op de fruitschaal lag. De banaan maakt het beslag
zoet en zorgt samen met het eiwit dat het bindt, dus er hoeft geen suiker of melk
bij.

\blockrule{Ingrediënten}
\begin{ingredients}
  \ing{3}{eiwitten}\index[register]{Eiwit}
  \ing{1}{overrijpe banaan}\index[register]{Banaan}
  \ing{30–40 g}{tarwebloem}\index[register]{Bloem}
  \ingb{ongezouten roomboter of zonnebloemolie, om in te bakken}
  \ingb{poedersuiker, stroop of hazelnootpasta, als topping}
\end{ingredients}

\blockrule{Bereiding}
\tip{Bak de pannenkoekjes klein. Zonder eidooier zijn ze wat kwetsbaarder en
breken grote exemplaren makkelijk bij het omdraaien.}
\begin{steps}
  \item Pel de banaan en breek hem in stukken in een hoge maatbeker. Voeg de
        eiwitten en de tarwebloem toe.
  \item Pureer met de staafmixer tot een glad beslag zonder klontjes.
  \item Verhit een koekenpan op middelhoog vuur met een klein beetje boter of
        olie, tot de boter schuimt of de olie glanzend en dun is.
  \item Schep kleine hoopjes beslag in de pan, met wat ruimte ertussen. Bak ze
        ongeveer 2 minuten tot de onderkant goudbruin is en de bovenkant begint
        te stollen, draai om met een spatel en bak de andere kant nog ongeveer
        2 minuten tot ze goudbruin en gaar zijn.
  \item Serveer met poedersuiker, stroop of hazelnootpasta. Hazelnootpasta met
        een paar plakjes extra banaan is ook lekker.
\end{steps}
\end{recipe}
